\section{Presented probability-preserving operators}\label{sec:operators}

\subsection{Finite kernels and their exact identities}

Let $(X,\mu)$ and $(Y,\nu)$ be finite probability spaces. A
probability-preserving transition is a rational matrix $P=(P_{ij})$ with
\begin{equation}\label{eq:stationarity}
 P_{ij}\ge0,\quad \sum_jP_{ij}=1,\quad
 \sum_i\mu_iP_{ij}=\nu_j.
\end{equation}
Its action on observables is contravariant:
\[
 T_P:L^1(Y,\nu)\longrightarrow L^1(X,\mu),\qquad
 (T_Pf)_i=\sum_jP_{ij}f_j.
\]
Zero-mass source rows may be chosen arbitrarily subject to row
normalization; they have no effect on the $L^1$ map. They must nevertheless
remain well-formed rows in the numerical code.

\begin{theorem}[Finite Markov-operator calculus]\label{thm:finite-markov}
The map $T_P$ is positive, linear, preserves $1$ and the integral, and
\begin{equation}\label{eq:markov-contraction}
 \norm{T_Pf-T_Pg}_{1,\mu}\le\norm{f-g}_{1,\nu}.
\end{equation}
If $Q:(Y,\nu)\to(Z,\xi)$ is another transition, its matrix composite $PQ$
satisfies \eqref{eq:stationarity}, and $T_{PQ}=T_PT_Q$.
Finite independent products are preserved by the tensor matrix
$(P\otimes Q)_{(i,k),(j,\ell)}=P_{ij}Q_{k\ell}$.
\end{theorem}
\begin{proof}
Positivity and linearity follow termwise; the row-sum law gives $T_P1=1$.
Finite Fubini and the last equation in \eqref{eq:stationarity} give
\[
 \sum_i\mu_i(T_Pf)_i
 =\sum_j\left(\sum_i\mu_iP_{ij}\right)f_j
 =\sum_j\nu_jf_j.
\]
For contraction, apply the triangle inequality before exchanging the sums:
\[
 \sum_i\mu_i\left|\sum_jP_{ij}(f_j-g_j)\right|
 \le\sum_{i,j}\mu_iP_{ij}|f_j-g_j|
 =\sum_j\nu_j|f_j-g_j|.
\]
For composition, positivity is termwise, the row sum is
$\sum_jP_{ij}\sum_kQ_{jk}=1$, and the new marginal is
$\sum_j\nu_jQ_{jk}=\xi_k$. Expanding a triple sum proves the action law
and associativity. The tensor equations factor into one equation for each
factor by finite Fubini. These proofs use only bounded sums and
nonnegative rational arithmetic.
\end{proof}

The existing \texttt{FiniteKernel} module concerns finite expectation
evaluation and natural-weight presentations; the matrix and completion
arguments here build on those results rather than changing their meaning
\cite{repo-finite-kernel,repo-fubini}. The general categorical setting is
related to Markov categories \cite{fritz}, with one structural restriction
explained below.

\subsection{Strong action data instead of uniform matrix approximation}

We now allow the finite probability algebras to refine indefinitely, as in
the measure presentation of the preceding section. Let $D_X,D_Y$ denote
their rational simple-function codes. Their $L^1$ completions are described
by regular names.

\begin{definition}[A presented bounded action]\label{def:action}
A bounded action $A:X\to Y$ on observables consists of a rational bound
$C\ge0$ and a PR routine $A(s,p)\in D_Y$ for $s\in D_X$. After common
refinement its finite contracts are
\begin{align}
 \norm{A(s,p)-A(t,q)}_1
 &\le C\norm{s-t}_1+2^{-p}+2^{-q},\label{eq:action-stability}\\
 \norm{A(as+bt,p)-aA(s,q)-bA(t,r)}_1
 &\le2^{-p}+|a|2^{-q}+|b|2^{-r},\label{eq:action-linearity}\\
 \norm{A(s,p)}_1&\le C\norm s_1+2^{-p}.\label{eq:action-bound}
\end{align}
They hold for all rational $a,b$, simple $s,t$, and natural precisions.
For a probability-preserving action require additionally $C=1$ and
\begin{align}
 \norm{(A(s,p))_-}_1&\le2^{-p}\quad(s\ge0),\label{eq:action-positive}\\
 \norm{A(1,p)-1}_1&\le2^{-p},\qquad
 |I_Y(A(s,p))-I_X(s)|\le2^{-p}.\label{eq:action-preserve}
\end{align}
Here $u_-=(|u|-u)/2$ is defined on simple functions and extends
nonexpansively to $L^1$. Each displayed inequality is a numerical
all-precision contract for a fixed routine.
\end{definition}

The definition does not assume a map on an already constructed uncountable
space. It gives its finite action and the precise relations needed to
construct one. For actual implementations the contracts must be proved;
placing them in a Lean structure is not a proof that an implementation
inhabits that structure.

\begin{theorem}[Extension and composition of strong actions]\label{thm:action-extension}
A presented bounded action extends uniquely to a bounded linear map on
regular $L^1$ names. Its bound is $C$. Probability-preserving actions extend
to positive unital integral-preserving contractions. Identity and
composition have PR action routines, and equality on dense simple
functions implies equality on all names.
\end{theorem}
\begin{proof}
Choose a natural $\ell$ with $C\le2^\ell$; one can take the bit length of
$\lceil\max(1,C)\rceil$. If $f$ has simple approximants $s_n$ with
$\norm{s_n-f}_1\le2^{-n}$, define
\[
 b_p=A(s_{p+\ell+2},p+2).
\]
For $p,q$ the stability contract gives
\[
 \norm{b_p-b_q}_1
 \le C(2^{-p-\ell-2}+2^{-q-\ell-2})
      +2^{-p-2}+2^{-q-2}
 \le2^{-p-1}+2^{-q-1}.
\]
Thus $b$ is a regular name. Comparing two equivalent input names with
\eqref{eq:action-stability} shows the output is independent of the name.
Taking arbitrarily precise finite comparisons gives
$\norm{Af-Ag}_1\le C\norm{f-g}_1$. The linearity contract gives exact
linearity by the same process. The negative-part, unit, and integral
contracts pass to the limit because these operations are nonexpansive in
$L^1$.

For actions $A:X\to Y$ and $B:Y\to Z$, an evaluator for $BA$ is
\[
 (B\circ A)(s,p)=B(A(s,p+\ell_B+2),p+2),
 \qquad C_B\le2^{\ell_B}.
\]
Its distance from the intended composite is at most
$C_B2^{-p-\ell_B-2}+2^{-p-2}\le2^{-p-1}$.
Triangle inequalities with the separate action contracts prove all the
required finite laws, with one extra guard bit if one insists on an exact
$2^{-p}$ convention. The bound is $C_BC_A$. The identity routine returns
its input exactly. The two associations of three actions agree on every
simple function, hence on every name by continuity. Finally, if two maps
agree on simple functions, comparison through a simple approximant $s_n$
gives output distance at most $(C_A+C_B)2^{-n}$; this proves uniqueness.
\end{proof}

\begin{example}[Why strong action is needed]\label{ex:no-operator-norm}
On the dyadic presentation of $[0,1]$, let $E_n$ average over the cells of
the $n$th partition. For every supplied $L^1$ name $f$, the previous
section constructs a PR rate for $E_nf\to f$. But for every $n$ there is
a simple function $h_n$ equal to $+1$ and $-1$ on the two halves of every
$n$th-level cell. Then $\norm{h_n}_1=1$ and $E_nh_n=0$, so
\[
 \norm{I-E_n}_{L^1\to L^1}\ge1.
\]
No sequence of these finite conditional averages approximates the identity
in operator norm. A definition that required this convergence would
exclude the identity from the intended completion. Strong action on each
named input avoids that error and provides precisely the data needed for
composition.
\end{example}

\subsection{Categorical structure and its probability constraint}

Take presented probability algebras as objects, and a transition $X\to Y$
to mean a probability-preserving action $L^1(Y)\to L^1(X)$ modulo equality
on names. Theorem~\ref{thm:action-extension} gives a category. The word
``category'' can be eliminated from any fixed argument: objects and arrows
are presentation codes, and the associativity and identity statements are
all-precision arithmetic equalities. There is no need for a universe of
all spaces inside iPRA.

For products use the rectangular finite algebra with masses $\mu_i\nu_j$.
The finite tensor matrices in Theorem~\ref{thm:finite-markov} are
contractions on this algebra. More generally, to tensor two presented
actions, write a rectangular simple function as a finite sum
$\sum_{i,j}c_{ij}\mathbf1_{A_i}\mathbf1_{B_j}$ and apply the actions
separately to the two indicators. For a product of named functions,
\[
 \norm{fg-\widetilde f\widetilde g}_{L^1(X\times Y)}
 \le\norm{f-\widetilde f}_{L^1(X)}\norm g_{L^1(Y)}
    +\norm{\widetilde f}_{L^1(X)}\norm{g-\widetilde g}_{L^1(Y)}.
\]
Finite Fubini proves this directly. It provides the PR precision schedule
for the finite tensor sum. Positivity, mass preservation, and the pointwise
inequality $|Tu|\le T|u|$ give contraction on rectangular simple
functions explicitly: put $a_i=T\mathbf1_{A_i}$ and
$b_j=S\mathbf1_{B_j}$. These functions are nonnegative, so
$|\sum c_{ij}a_ib_j|\le\sum|c_{ij}|a_ib_j$; integrating gives
$\sum|c_{ij}|\mu(A_i)\nu(B_j)$, the norm of the original atom table.
For a finite evaluator, clip indicator-image approximants to $[0,1]$.
If each has $L^1$ error at most $\delta$, each product has error at most
$2\delta$. The choice
$\delta\le2^{-p}/(2\max(1,\sum_{ij}|c_{ij}|))$ suffices for the whole
table, and only one image per source atom needs to be cached.
The tensor therefore extends to the product $L^1$ completion. Common
refinement and linearity prove independence of the chosen rectangular
decomposition. The tensor unit is the one-point probability space;
symmetry and associativity are the explicit index permutations.

This category has designated marginal measures. Copying $x\mapsto(x,x)$
generally sends $\mu$ to a diagonal law, not to $\mu\otimes\mu$.
For a fair two-point space the diagonal has mass $1$ under the copied law
and mass $1/2$ under the product law. Thus copying is not in general an
arrow to this tensor product. The full category of Markov kernels studied
in \cite{fritz} does carry copy maps; the probability-preserving category
here has the stated restricted structure. Confusing these two categories
would give incorrect independence arguments. We assert only the structure
just proved.

\subsection{Kernel perturbations and conditioning}

\begin{proposition}[A tail-sensitive kernel bound]\label{prop:kernel-tail}
Suppose $P,Q$ have the same marginals $\mu,\nu$, and set
\[
 d_\mu(P,Q)=\sum_i\mu_i\sum_j|P_{ij}-Q_{ij}|.
\]
For bounded $f$, $|f|\le K$, one has
$\norm{T_Pf-T_Qf}_{1,\mu}\le Kd_\mu(P,Q)$.
For arbitrary integrable $f$ and any $K>0$,
\begin{equation}\label{eq:kernel-tail}
 \norm{T_Pf-T_Qf}_{1,\mu}
 \le2t_K(f)+Kd_\mu(P,Q).
\end{equation}
\end{proposition}
\begin{proof}
For bounded $f$, expand the difference of the two rows and bound
$|f_j|$ by $K$ before summing. For general $f$, insert $c_K(f)$ and use
the two $L^1$ contraction bounds on $f-c_K(f)$. Each of these costs
$t_K(f)$; the remaining bounded difference costs $Kd_\mu(P,Q)$.
Completion is valid because clipping and all the operators are continuous
in $L^1$. Hence a supplied tail modulus and a kernel-approximation rate
give a strong-action rate even when no useful uniform bound on $f$ exists.
\end{proof}

\begin{proposition}[Conditioning with a stated mass margin]\label{prop:bayes}
Let $a_j,b_j\ge0$ have masses $A=\sum_ja_j$, $B=\sum_jb_j$ with
$A,B\ge\tau>0$. Put $\delta=\sum_j|a_j-b_j|$. Then
\begin{equation}\label{eq:bayes}
 \sum_j\left|\frac{a_j}{A}-\frac{b_j}{B}\right|
 \le\frac{2\delta}{\tau}.
\end{equation}
Thus bounded posterior expectations have a PR modulus once the positive
mass margin is part of the input.
\end{proposition}
\begin{proof}
Add and subtract $b_j/A$, sum, and use $\sum_jb_j=B$:
\[
 \sum_j\left|\frac{a_j}{A}-\frac{b_j}{B}\right|
 \le\frac{\delta}{A}+B\left|\frac1A-\frac1B\right|
 =\frac{\delta+|A-B|}{A}\le\frac{2\delta}{\tau}.
\]
The inequality $|A-B|\le\delta$ is finite weighted error. Clearing the
positive denominators gives a quantifier-free arithmetic proof.
\end{proof}

The margin is necessary for a uniform pointwise posterior. The unnormalized
rows $(\eps,0)$ and $(0,\eps)$ are distance $2\eps$ apart, while their
normalized laws are distance $2$ apart. Conditional expectations as $L^1$
classes instead weight each posterior by its cell mass; the cancellation
of that denominator is exactly what makes the tower construction stable.
For named real masses we do not decide whether an arbitrary mass is zero.
Finite rational masses can be tested, and continuous posterior queries
must carry a positive margin. General disintegration is an additional
construction, not a consequence of the finite tower identity alone.
